\section{Lean formalization and its remaining steps}\label{sec:lean}

The companion \texttt{lean/Brouwer.lean} uses Lean 4.19.0 and its standard
library. Its proofs concern executable Lean definitions and integer
inequalities. ProveIt's embedded judgment is
\texttt{Prf PRAAx} in \texttt{LeanProofs.ConciseIPRA}. The pinned repository
uses Lean 4.32.0. A proof in the former setting supplies mathematics for the
reification work in the latter. This packet includes a kernel-checked
certificate-to-error theorem and a human proof of the general parity step.

\subsection{What the companion checks}
After clearing denominators, a vertex and a sample are integer rows having
the same positive total mass. For instance choose a common multiple $J$ of
$N$, all coordinate denominators of the selected sample rows and $y$, and the
denominator of $\omega$. Multiplication by $J$ preserves the color comparisons.
The integer tolerances for the local theorem are $J/N$ and $J\omega$.

\begin{center}\small
\begin{tabularx}{\linewidth}{@{}>{\raggedright\arraybackslash}p{0.36\linewidth}>{\raggedright\arraybackslash}X@{}}
\toprule
\textsf{Lean declaration} & \textsf{Mathematical content}\\
\midrule
\texttt{maxLabel\_correct} &
A label exists on positive support and attains the global residual maximum.\\[2pt]
\texttt{stop\_mesh}, \texttt{gap\_mesh} &
Cumulative coordinates are monotone along a prefix. Each gap changes by at most one.\\[2pt]
\texttt{solveWord\_sound}, \texttt{solveWord\_complete} &
The bounded search returns a full word when one exists, and every returned word is in range and passes the full-color test.\\[2pt]
\texttt{certified\_error}, \texttt{solveWord\_error} &
The full-color bit supplies each color witness. A successful grid search and the local sample bounds imply the $2d/(d+1)$ estimate.\\[2pt]
\texttt{cubeLabel\_correct}, \texttt{cube\_anchor} &
The cubical labeling satisfies all three boundary requirements. Its two residual signs give the cubical error bound.\\
\bottomrule
\end{tabularx}
\end{center}
In particular, \texttt{solveWord\_complete} takes existence of a full word
as a premise. Section~\ref{sec:finite-sperner} discharges that premise
mathematically. An implementation of finite search alone would leave precisely
this point open.

\subsection{Completing the embedded derivation}
The following modules can be developed in dependency order. All formula
parameters are numerical variables or fixed PR terms as in
Section~\ref{sec:formal}.
\begin{enumerate}
\item \emph{Arithmetic and tables.} Introduce PR codes for indexing,
word digits, canonical rationals, table assembly, the full-color bit, parity
and bounded least success. Instantiate the repository's
\texttt{Core.Basic}, \texttt{Core.BoundedSums}, \texttt{Core.SumOrder} and
\texttt{Analysis.RationalArithmetic}. Prove lookup after update by induction
on the index. Prove the prefix-count and least-hit invariants given above.
A total signed integer can instead be represented as a pair of naturals. Here
$(p,n)$ denotes $p-n$, and $(p,n)\le(q,m)$ abbreviates $p+m\le q+n$.
Thus all scalar comparisons used in the transfer lemmas can be expressed in
the natural-number object language.
\item \emph{Uniform grid parity.} Code a cell by the word of length $r$
and an incidence by $(r+1)\,\mathrm{word}+k$. Implement the permutation
by sorting the time positions by (digit ascending, position descending).
The inverse sends a valid permutation to its sequence of base entries.
Prove the tie-block characterization by checking adjacent entries; verify
the inverse by counting the earlier appearances of each value. For each of
the three deletion cases, encode the partner from
Section~\ref{sec:finite-sperner}, with a sentinel for the unpaired cases.
Prove, by bounded vertex comparison, preservation of the retained list,
the inverse equation, and the positive-gap exclusions. Express the
top-face bijection as append-and-drop identities on the word/base
representation. The resulting finite double sums give the one-step
parity equality. Apply object-level quantifier-free induction to the
capped dimension term with the ambient table held fixed. Use
odd-count-implies-first-hit to finish the search proof.
\item \emph{Rational extraction.} Reify the inequalities in
\texttt{certified\_error}. For a color $i$, obtain its first occurrence
using the Boolean full-color proof. Compare each $j$ to $i$ at that
occurrence, transfer to the anchor and sum. A running sum with its
diagonal term omitted supplies the multiplier $d$. Reify
\eqref{eq:main} by substituting $h=1/N$, $\omega=2^{-m}$ and
cross-multiplying. Introduce arbitrary $s\ge m$ after building the
table and generalize it after the estimate. Close the universal formula
of Section~\ref{sec:formal} by the object rules.
\item \emph{Shorter fixed-dimension route.} For the cube instantiate
\texttt{sperner\_cube} in the namespace\linebreak[1] \texttt{LeanProofs.IPRA.Topology.SpernerCube}
with the term for the least successful cubical coordinate.
The evaluator parameters and the desired precision can be packed into
the provider's parameter $p$. The coordinate $j$ of the mathematical
labeling is at list depth $d-1-j$ in the provider. Reify the three
parts of \texttt{cubeLabel\_correct} for the canonical grid codes.
The provider then yields $\operatorname{par}(\mathrm{count})=1$.
Scan the provider's bases and permutations for a full cell, use the two
sign inequalities in \texttt{cube\_anchor}, and instantiate the mixed
modulus to obtain a cube witness. For the ball, add $R$, its two
polynomial distance identities after clearing denominators, and the affine
maps $I,P$. One can also obtain a fixed-dimensional simplex result by applying the cube theorem to $A\circ\Pi$ in $d+1$ variables and returning $\Pi(u)$. The inclusion of the simplex in the cube is isometric and $\Pi$ is a $2$-Lipschitz retraction.
\end{enumerate}

The upstream fixed-dimension parity source was inspected at the pinned
revision. The full object-language derivation of approximate Brouwer,
including the dimension-uniform parity derivation, remains to be implemented.
The supplied Lean command and Python checks can be run independently as
described in \texttt{README.md}.
